% !TeX program = xelatex
% Overleaf 编译器请选择 XeLaTeX。
\documentclass[a4paper,8pt]{article}

\usepackage[margin=1.25cm]{geometry}
\usepackage[UTF8,fontset=fandol]{ctex}
\usepackage{fontspec}
\usepackage{enumitem}
\usepackage{hyperref}
\usepackage{xcolor}
\usepackage{titlesec}
\usepackage{tabularx}
\usepackage{array}
\usepackage{microtype}
\usepackage{graphicx}

% =========================================================
% 字体
% =========================================================

\setmainfont{lmroman10-regular.otf}[
    BoldFont = lmroman10-bold.otf,
    ItalicFont = lmroman10-italic.otf,
    BoldItalicFont = lmroman10-bolditalic.otf,
    BoldFeatures = {FakeBold=1.8}
]

\setCJKmainfont{FandolSong-Regular.otf}[
    BoldFont = FandolSong-Bold.otf,
    BoldFeatures = {FakeBold=2.2}
]

% =========================================================
% 颜色
% =========================================================

\definecolor{main}{HTML}{006F73}
\definecolor{textgray}{HTML}{111111}
\definecolor{muted}{HTML}{444444}

\hypersetup{
    colorlinks=true,
    urlcolor=main,
    linkcolor=main
}

\urlstyle{same}

\pagestyle{empty}
\color{textgray}

% =========================================================
% 全局排版
% =========================================================

\setlength{\parindent}{0pt}
\setlength{\parskip}{0pt}
\setlength{\tabcolsep}{0pt}

% 正文行距统一；为新增开源贡献保持单页而略作收紧
\renewcommand{\baselinestretch}{1.10}

% 避免英文过多时行内空格被拉得太夸张
\emergencystretch=1em

% 中文换行更自然
\XeTeXlinebreaklocale "zh"
\XeTeXlinebreakskip = 0pt plus 1pt

% =========================================================
% Bullet 列表
% =========================================================

\setlist[itemize]{
    leftmargin=1.25em,
    itemsep=2.6pt,
    topsep=2.5pt,
    parsep=0pt,
    partopsep=0pt,
    label=\textcolor{main}{\small\textbullet}
}

% =========================================================
% Section
% =========================================================

\titleformat{\section}
{\large\bfseries\color{main}}
{}{0pt}{}
[\vspace{-1pt}\titlerule]

% 左、上、下
\titlespacing*{\section}
{0pt}
{6pt}
{4pt}

\newcommand{\resumeItem}[1]{\item #1}

\newcommand{\skill}[2]{%
    \textbf{#1：}#2\par\vspace{1.5pt}
}

% =========================================================
% 实习标题
%
% 左：公司
% 中：岗位
% 右：时间
% =========================================================

\newcommand{\experienceTitle}[3]{%
    \begin{tabularx}{\textwidth}{
        @{}
        >{\raggedright\arraybackslash}p{0.40\textwidth}
        >{\centering\arraybackslash}p{0.36\textwidth}
        >{\raggedleft\arraybackslash}p{0.24\textwidth}
        @{}
    }
        \textbf{\large #1}
        &
        \textbf{#2}
        &
        \textcolor{muted}{#3}
    \end{tabularx}
}

% =========================================================
% 项目标题
%
% 左：项目名称
% 右：时间
% 下一行技术栈
% 下一行项目链接
% =========================================================

\newcommand{\projectTitle}[4]{%
    \begin{tabularx}{\textwidth}{
        @{}
        >{\raggedright\arraybackslash}X
        >{\raggedleft\arraybackslash}p{0.18\textwidth}
        @{}
    }
        \textbf{\large #1}
        &
        \textcolor{muted}{#2}
    \end{tabularx}

    \vspace{1.5pt}

    \textbf{技术栈：}#3\\[1pt]
    \textbf{项目链接：}\href{#4}{\nolinkurl{#4}}
}

\begin{document}

% =========================================================
% 个人信息
% =========================================================

\noindent
\begin{minipage}[t]{0.43\textwidth}
    \vspace{0pt}

    {\fontsize{22pt}{24pt}\selectfont\bfseries 孔令航}

    \vspace{10pt}

    {\large\bfseries
        可实习 6 个月及以上
    }

    \vspace{2pt}

    {\normalsize\bfseries
        两周内到岗
    }

\end{minipage}
\hfill
\begin{minipage}[t]{0.37\textwidth}
    \vspace{2pt}

    \raggedleft
    \small

    电话：15832895339\\[3pt]

    邮箱：
    \href{mailto:17736236765@163.com}
    {17736236765@163.com}\\[3pt]

    GitHub：
    \href{https://github.com/Kong-lh-rgb}
    {github.com/Kong-lh-rgb}

\end{minipage}
\hfill
\begin{minipage}[t]{0.14\textwidth}
    \vspace{0pt}

    \raggedleft
    \includegraphics[
        width=2.05cm,
        keepaspectratio
    ]{resume-photo-blue.png}

\end{minipage}

\vspace{-11pt}

% =========================================================
% 教育背景
% =========================================================

\section{教育背景}

\begin{tabularx}{\textwidth}{
    @{}
    >{\raggedright\arraybackslash}X
    >{\raggedright\arraybackslash}X
    >{\centering\arraybackslash}p{0.10\textwidth}
    >{\raggedleft\arraybackslash}p{0.18\textwidth}
    @{}
}
    \textbf{内蒙古大学（211）}
    & 计算机技术
    & 硕士
    & 2025--2028
    \\[1pt]

    \textbf{河北金融学院}
    & 计算机科学与技术
    & 本科
    & 2021--2025
\end{tabularx}

% =========================================================
% 实习经历
% =========================================================

\section{实习经历}

\experienceTitle
{杭州广立微电子股份有限公司}
{AI Agent 评测实习生}
{2026.04--2026.07}

\vspace{-8pt}

参与内部 AI Agent 自动化评测平台 QAtools 建设，
负责异构 Agent 平台接入、批量回归评测及长任务执行链路优化。

\begin{itemize}

\resumeItem{
    \textbf{Agent 平台接入：}
    抽象统一 AgentAdapter 协议，完成 SemiMind、SemiClaw 等异构 Agent 平台接入，统一配置管理、Task 创建、WebSocket 流式事件与结果模型，支撑内部 Agent 的批量自动化回归评测。
}

\resumeItem{
    \textbf{评测调度与性能优化：}
    设计 Dataset $\rightarrow$ Agent Run $\rightarrow$
    Result Parse $\rightarrow$ Evaluate 批量评测 Workflow，
    基于异步 Worker Pool、并发限流及可重试异常处理优化执行效率，
    将 \textbf{1,000 条 Case 单轮评测耗时由约 3.2 h 降至 50 min，
    降低约 74\%}。
}

\resumeItem{
    \textbf{长任务执行可靠性：}
    重构 Agent Runner，
    将 WebSocket 连接状态与远端 Task 生命周期解耦，
    基于 \texttt{task\_id} 实现断线重连、状态补查、
    结果补拉与事件幂等去重，降低批量测试中因流式连接异常导致的任务误失败与重复执行。
}

\end{itemize}

% =========================================================
% 项目经历
% =========================================================

\section{项目经历}

\projectTitle
{Vesta｜面向长程任务的本地 AI Agent Harness}
{2026.07--至今}
{Python、FastAPI、SQLite、MCP、Electron、React、TypeScript、Swift、OpenAI SDK}
{https://github.com/Kong-lh-rgb/vesta}

\vspace{1pt}

独立从 0 到 1 设计并实现本地 AI Agent Harness，
打通多模型适配、Agent Loop、Tool Calling、长任务执行、
上下文治理、长期记忆、安全沙箱、Checkpoint/Recovery 与 Desktop 完整链路。

\begin{itemize}

\resumeItem{
    \textbf{Agent Runtime 与工具执行：}
    设计并实现统一 \textbf{Agent Runtime}，
    抽象 Model Adapter、Agent Loop 与
    \textbf{Registry--Executor--Hook} 工具执行层，
    统一编排内置工具、MCP、Skill 与 Computer Use；
    建立参数校验、超时、审批等异常处理机制，
    并通过 WebSocket + JSON-RPC 打通 Electron、Python Host 与 Swift Helper，
    实现从模型决策到 macOS 原生操作的完整执行链路。
}

\resumeItem{
    \textbf{长程任务与故障恢复：}
    面向 Long-Horizon Agent 设计
    \textbf{Conversation--Run--Task} 分层状态模型，
    将任务状态与单次进程生命周期解耦；
    基于 Checkpoint 与 Recovery Run 保存关键执行状态，
    支持模型异常、用户中断及 Host 重启后的任务恢复，
    避免长程任务失败后整体重新执行。
}

\resumeItem{
    \textbf{上下文与 Token 优化：}
    将平均可计费 Token 从
    \textbf{3,766 降至 2,767（-26.5\%）}，
    Prompt Cache 平均命中率由
    \textbf{68.0\% 提升至 75.5\%}；
    通过解耦完整会话历史与模型实际上下文，
    设计分层裁剪、滚动压缩、稳定前缀及 Run 级 Token Budget，
    在保持长程上下文可用性的同时降低模型调用成本。
}

\resumeItem{
    \textbf{长期记忆与经验学习：}
    构建 \textbf{Core/Ordinary Memory} 分层长期记忆，
    形成“\textbf{自动召回--按需检索--反思更新}”闭环；
    结合 \textbf{FTS5 与向量语义检索并通过 RRF 融合排序}，
    以 Post-Run Reflection 和 Revision 实现跨会话记忆更新，
    并从 Completed Task 与真实 Trace 中提炼Skill，形成可复用 Agent 经验。
}

\resumeItem{
    \textbf{Agent Eval 与可靠性：}
    构建端到端 \textbf{Eval Harness}，
    自动记录 Trace、中间状态、失败原因与版本化 Baseline，
    建立\textbf{评测--失败定位--优化--回归}闭环；
    针对 68 个能力单元使用 DeepSeek V4 Flash 各执行 3 次，
    完成 \textbf{204 次真实模型回归}，
    达到 \textbf{192/204 样本通过}，
    其中 \textbf{57/68 能力单元连续三次稳定通过}。
}

\end{itemize}

% =========================================================
% 开源贡献
% =========================================================

\section{开源贡献}

\begin{tabularx}{\textwidth}{@{}X>{\raggedleft\arraybackslash}p{0.18\textwidth}@{}}
    \textbf{\large
        \href{https://github.com/multica-ai/multica}
        {Multica｜开源 AI Agent 协作平台（48.8k+ Stars）}}
    &
    \textcolor{muted}{2026.09}
\end{tabularx}

\vspace{1pt}

\textbf{核心代码贡献：}
修复共享浅克隆缓存遗漏最新提交、导致 Agent 隔离工作目录检出失败的问题；
调整 Repo Cache 引用同步与 Checkout 顺序，构造确定性回归测试，
完整后端 CI 通过，\href{https://github.com/multica-ai/multica/pull/8024}{PR \#8024} 已由维护者合并。

% =========================================================
% 专业技能
% =========================================================

\section{专业技能}

\small

\skill{编程与后端工程}{
    Python、Go、SQL；
    FastAPI、Redis；
    PostgreSQL、Docker、Linux；
    熟悉异步编程、并发控制、流式通信及常见后端服务开发
}

\skill{LLM 与 Agent 系统}{
    Structured Output、Function Calling、
    ReAct、MCP、Skills、Computer Use；
    Agent Runtime/Harness、熟悉 LangGraph 任务编排、上下文管理、
    长期记忆、Human-in-the-Loop、Checkpoint/Recovery
}

\skill{Agent 工程与可靠性}{
    多模型接入、长程任务执行、工具执行与状态管理；
    Prompt Cache、Token Budget、超时/重试、故障恢复；
    Trace/日志分析、Agent Eval、
    LLM-as-a-Judge、规则与状态断言
}

\skill{模型与知识检索}{
    Transformer、PyTorch、Hugging Face Transformers；
    RAG、Embedding、BGE-M3、Milvus、Hybrid Search、
    KG-RAG
}

\end{document}
